\section{结论}\label{sec:conclusion}

本文用两个任务回答同一个问题：图算法里的“白盒”与“黑盒”，到底该怎么选。

\begin{enumerate}[leftmargin=1.8em]
\item \textbf{算法已知时，结构就是算法。} $k$ 层 min-sum 消息传递与同步 Bellman-Ford 的第 $k$ 轮逐点等价（定理~\ref{thm:equiv}），$L$ 层模型的感受野就是 $L$ 跳，因此深度不足是结构性错误而非训练不足（命题~\ref{prop:depth}）。
在实验上，未训练的恒等 min-sum 误差 $\sim10^{-6}$，训练 2 个参数反而退化到 $0.024$——精确解本来就在假设类里，学习只会把它推开。
白盒的鲁棒性不是“更抗打”，而是三条可证命题（命题~\ref{prop:mono}--\ref{prop:bound}）：删边单调与非树边不变各 500/500 零违反，单边删除上界在 398 个非平凡例上零违反。
\item \textbf{算法未知时，问题变成“能不能验证解释”。} Cora 上 GCN 0.8162 对 MLP 0.5539（参数量相同），FGA 攻击比随机强 7 倍；
但真实数据没有“正确的解释”，只能造一个带真值的 SBM 环境，测得解释器达到 oracle 的 93.9\%（随机边 6\%）。
而攻击把准确率打掉 7.5 个点、解释却几乎不变（Jaccard 1.000），说明\textbf{解释可以稳定，却不跟踪模型改变的主意}。
\item \textbf{方法论上，先问指标会不会被骗，再问数值是多少。} 第一部分用常数预测器骗过了 Fidelity$^+$（28 图集上 0.978 配 0.13，60 图集上 $-0.10$）；第二部分用 bottom-$k$ 负控制证明主判据有“牙”，并指出充分性/必要性/稳定性三者都只是必要条件。
\end{enumerate}

合起来：\textbf{“可解释性 vs. 性能”的权衡曲线只在算法未知时才真正出现}。算法已知时，把算法写进结构能同时赢下三轴；算法未知时，取舍存在，但“可解释”本身先要有一个可证伪的定义——这个代价往往被低估。

\subsection*{局限与后续工作}

本文的局限已在第~\ref{sec:limits}~节逐条列出，最重要的三条是：第一部分的黑盒基线没有调到最好（6/8 个种子塌缩）；两部分的参数量与迭代预算都不对称；第二部分的“真值”来自合成图，不真实。

后续工作有三个方向：
（1）把黑盒换成权重共享的结构，让它也能迭代到收敛，消掉迭代预算这个不对称；
（2）设计\emph{对模型改变敏感}的解释器，检验“攻击前后解释不变”这一现象是否可以被打破；
（3）把任务从“有精确解”扩展到“有近似算法”（如旅行商），看权衡曲线是否以另一种形式出现。

\subsection*{复现}

全部实验在纯 CPU 上完成。第一部分依赖 \texttt{numpy}、\texttt{networkx}、\texttt{matplotlib}、\texttt{torch}，不含外部数据集；第二部分额外依赖 \texttt{scipy}，Cora 由 \texttt{python exp/fetch\_cora.py} 下载（约 570 KB，不入库）。
仓库结构、运行命令与结果文件见 \texttt{README.md}；每个实验都有对应的验收脚本，并配有“桩实现必须失败”的负控制（第一部分 9 条，第二部分 3 条）。
